%\usepackage[utf8]{inputenc} % Umlaute, other option latin1
%\usepackage{amsmath}
%\usepackage{amssymb}
\usepackage{booktabs} % toprule, midrule, bottomrule
\usepackage{adjustbox} % frame in includegraphics [export]

% tables
\usepackage{array} % tabular m{1cm}
\usepackage{ragged2e} % own column types P{1cm}
\newcolumntype{P}[1]{>{\RaggedRight\hspace{0pt}}p{#1}}
\newcolumntype{M}[1]{>{\RaggedRight\hspace{0pt}}m{#1}}

% TODOs
\newcommand{\todo}[1]{{\color{red} [#1]}}
\newcommand{\todoo}[1]{{\color{blue} \{#1\}}}
\newcommand{\todooo}[1]{{\color{green} (#1)}}
\newcommand{\changed}[1]{{\color{blue} #1}}
\newcommand{\changedd}[1]{{\color{blue} #1}}
\newcommand{\changeddd}[1]{{\color{blue} #1}}
%\newcommand{\changed}[1]{{\textbf{#1}}}
\newcommand{\fodo}[1]{\todo{\footnote{\todo{#1}}}}
\newcommand{\fodoo}[1]{\todoo{\footnote{\todoo{#1}}}}
\newcommand{\fodooo}[1]{\todooo{\footnote{\todooo{#1}}}}
\newcommand{\tocite}[1]{{\color{red} [cite:#1]}}
\newcommand{\todots}{\todo{\ldots}}
\newcommand{\kword}[1]{{\LARGE #1}}

%\renewcommand{\todo}[1]{}
%\renewcommand{\todots}{}
%\renewcommand{\tocite}[1]{}

%\renewcommand{\todoo}[1]{}
%\renewcommand{\todooo}[1]{}
%\renewcommand{\changed}[1]{#1}
%\renewcommand{\changedd}[1]{#1}
%\renewcommand{\changeddd}[1]{#1}
%\renewcommand{\fodo}[1]{}
%\renewcommand{\fodoo}[1]{}
%\renewcommand{\kword}[1]{#1}

\newcommand{\feature}[1]{\emph{#1}}
\newcommand{\spl}[1]{\emph{#1}}
\newcommand{\code}[1]{\mbox{\sf#1}}
\newcommand{\tool}[1]{\textsc{#1}}

\newcommand{\KeY}{Ke\kern-0.1emY}
\newcommand{\MonKeY}{Mon\KeY}
%\newcommand{\JPF}{\tool{Java PathFinder}}
\newcommand{\JPF}{\tool{JPF}}

% Email-Adressen
%\newcommand{\email}[1]{
	%\href{mailto:#1}{#1}
%}

% Aussagenlogische Ausdr�cke
\newcommand{\pand}{\wedge}
\newcommand{\por}{\vee}
\newcommand{\pnot}{\neg}
\newcommand{\pequals}{\Leftrightarrow}
\newcommand{\pimplies}{\Rightarrow}
\newcommand{\pnimplies}{\nRightarrow}
\newcommand{\patmostone}{\mbox{\textit{atmost1}}}
\newcommand{\pchooseone}{\mbox{\textit{choose1}}}
\newcommand{\ptrue}{\mbox{\textit{true}}}
\newcommand{\pfalse}{\mbox{\textit{false}}}
\newcommand{\true}{\texttt{true}}
\newcommand{\false}{\texttt{false}}

% Richtig, falsch
\newcommand{\richtig}{
	\begin{tikzpicture}[yscale=.2,xscale=.2,overlay]
		\draw[line width=4,line join=bevel,green] (0,1) to [bend left=30] (1,0) to [bend left=15] (3,3);
	\end{tikzpicture}
}
\newcommand{\folgerichtig}{
	\begin{tikzpicture}[yscale=.2,xscale=.2,overlay]
		\draw[line width=4,line join=bevel,blue] (0,1) to [bend left=30] (1,0) to [bend left=15] (3,3);
	\end{tikzpicture}
}
\newcommand{\falsch}{
	\begin{tikzpicture}[yscale=.2,xscale=.2,overlay]
		\draw[line width=4,line join=bevel,red] (0,2) to [bend left=5] (2,0) (0,0) to [bend left=15] (2,2);
	\end{tikzpicture}
}
\newcommand{\fsel}{
	\begin{tikzpicture}[yscale=.15,xscale=.15,overlay]
		\draw[line width=4,line join=bevel] (0,1) to [bend left=30] (1,0) to [bend left=15] (3,3);
		\draw[line width=3,line join=bevel,green] (0,1) to [bend left=30] (1,0) to [bend left=15] (3,3);
	\end{tikzpicture}
}
\newcommand{\felim}{
	\begin{tikzpicture}[yscale=.15,xscale=.15,overlay]
		\draw[line width=4,line join=bevel] (0,2) to [bend left=5] (2,0) (0,0) to [bend left=15] (2,2);
		\draw[line width=3,line join=bevel,red] (0,2) to [bend left=5] (2,0) (0,0) to [bend left=15] (2,2);
	\end{tikzpicture}
}

%\newtheorem{theorem}{Theorem}[section]
%\newtheorem{lemma}[theorem]{Lemma}
%\newtheorem{proposition}[theorem]{Proposition}
%\newtheorem{corollary}[theorem]{Corollary}
%
%\newenvironment{proof}[1][Proof]{\begin{trivlist}
%\item[\hskip \labelsep {\bfseries #1}]}{\end{trivlist}}
%\newenvironment{definition}[1][Definition]{\begin{trivlist}
%\item[\hskip \labelsep {\bfseries #1}]}{\end{trivlist}}
%\newenvironment{example}[1][Example]{\begin{trivlist}
%\item[\hskip \labelsep {\bfseries #1}]}{\hfill$\square$\end{trivlist}}
%%\newenvironment{example}[1][Example]{\itshape\begin{trivlist}
%%\item[\hskip \labelsep {\bfseries #1}]}{\hfill$\square$\end{trivlist}\normalfont}
%\newenvironment{remark}[1][Remark]{\begin{trivlist}
%\item[\hskip \labelsep {\bfseries #1}]}{\end{trivlist}}
%
%\newcommand{\qed}{\nobreak \ifvmode \relax \else
      %\ifdim\lastskip<1.5em \hskip-\lastskip
      %\hskip1.5em plus0em minus0.5em \fi \nobreak
      %\vrule height0.75em width0.5em depth0.25em\fi}
      
\newcommand{\picsheight}{60mm}
\newcommand{\paperreference}[3]{
	\vspace{2mm}
	\tikz{\node[draw=gray,line width=.5pt,fill=white,inner sep=0mm,blur shadow={shadow blur steps=5,shadow blur extra rounding=2pt}] {\href{#3}{\includegraphics[height=\picsheight]{papers/#1}}};\node[yshift=-15mm,rounded corners,fill=uulmaccent,blur shadow={shadow blur steps=5,shadow blur extra rounding=2pt}] {#2};}
}
\newcommand{\authorsays}[3]{
	\tikz{\node[draw=gray,line width=1pt,fill=white,inner sep=0mm,blur shadow={shadow blur steps=5,shadow blur extra rounding=2pt}] {\includegraphics[height=\picsheight]{authors/#1}};\node[yshift=-25mm,rounded corners,fill=white,blur shadow={shadow blur steps=5,shadow blur extra rounding=2pt}] {#2 says ``#3''};}
}
